\documentclass[11pt,a4paper]{article}
% =====================================================================
%  Shared preamble for the Part V lecture notes (lectures 19-22).
%
%  These four lectures sit outside Geron, so the notes ARE the primary
%  source and are examined on the same terms as the textbook parts.
%  Everything here is chosen to make that readable on paper: one column,
%  generous measure, and the same six-colour palette as the slides so a
%  student moving between the two is not re-learning what red means.
%
%  Build with pdflatex (twice, for the toc and the refs):
%      cd notes && latexmk -pdf lecture-19.tex
%  or  make -C notes
% =====================================================================
\usepackage[T1]{fontenc}
\usepackage[utf8]{inputenc}
\usepackage{newtxtext,newtxmath}      % Times-like: prints well, reads well
\usepackage[scaled=0.86]{inconsolata} % code, at a size that sits beside Times
\usepackage{microtype}
% newtx defines \Bbbk, \openbox, \proof and \endproof; amssymb and amsthm
% then redefine them and abort. Freeing the four names before those two
% packages load is the standard fix, and it costs nothing here: we use
% amsthm only for \newtheorem.
\let\Bbbk\relax
\usepackage{amsmath,amssymb}
\let\openbox\relax
\let\proof\relax
\let\endproof\relax
\usepackage{amsthm}
\usepackage{booktabs}
\usepackage{enumitem}
\usepackage{xcolor}
\usepackage{tcolorbox}
\usepackage{titlesec}
\usepackage{graphicx}
\usepackage[hidelinks]{hyperref}
\tcbuselibrary{skins,breakable}

% ---- the course palette, from assets/css/custom.css -------------------
\definecolor{aimlprimary}{HTML}{0B3D62}   % deep blue  - structure
\definecolor{aimlaccent} {HTML}{C0392B}   % brick red  - failures
\definecolor{aimlsuccess}{HTML}{14663A}   % green      - repairs
\definecolor{aimlmath}   {HTML}{6C3483}   % purple     - the derivation
\definecolor{aimlmuted}  {HTML}{4B5563}
\definecolor{aimlrule}   {HTML}{B0BCC7}
\definecolor{aimlsoft}   {HTML}{F4F7F9}

\usepackage[a4paper,margin=32mm,top=30mm,bottom=30mm]{geometry}
\setlength{\parskip}{0.45\baselineskip}
\setlength{\parindent}{0pt}
\linespread{1.06}

% ---- headings --------------------------------------------------------
\titleformat{\section}{\Large\bfseries\color{aimlprimary}}{\thesection}{0.7em}{}
\titleformat{\subsection}{\large\bfseries\color{aimlprimary}}{\thesubsection}{0.7em}{}
\titleformat{\subsubsection}{\bfseries\color{aimlmuted}}{}{0em}{}
\titlespacing*{\section}{0pt}{1.6\baselineskip}{0.5\baselineskip}
\titlespacing*{\subsection}{0pt}{1.2\baselineskip}{0.35\baselineskip}

% ---- boxes -----------------------------------------------------------
% One box per role, and the role is the colour. A student who has seen the
% slides already knows what each one means before reading a word of it.
\newtcolorbox{keybox}[1][]{
  breakable, enhanced, colback=aimlsoft, colframe=aimlprimary,
  boxrule=0.4pt, left=8pt, right=8pt, top=6pt, bottom=6pt,
  fonttitle=\bfseries\small, coltitle=white, colbacktitle=aimlprimary,
  attach boxed title to top left={yshift=-2mm, xshift=4mm},
  boxed title style={boxrule=0pt, arc=1pt}, #1}

\newtcolorbox{failbox}[1][]{
  breakable, enhanced, colback=white, colframe=aimlaccent,
  boxrule=0.9pt, left=8pt, right=8pt, top=6pt, bottom=6pt,
  fonttitle=\bfseries\small, coltitle=white, colbacktitle=aimlaccent,
  attach boxed title to top left={yshift=-2mm, xshift=4mm},
  boxed title style={boxrule=0pt, arc=1pt}, #1}

\newtcolorbox{derivbox}[1][]{
  breakable, enhanced, colback=white, colframe=aimlmath,
  boxrule=0.6pt, left=8pt, right=8pt, top=6pt, bottom=6pt,
  fonttitle=\bfseries\small, coltitle=white, colbacktitle=aimlmath,
  attach boxed title to top left={yshift=-2mm, xshift=4mm},
  boxed title style={boxrule=0pt, arc=1pt}, #1}

% An exercise. Deliberately the same shape as the notebook's `try` field:
% one change, and what should happen to the output.
\newtcolorbox{trybox}[1][]{
  breakable, enhanced, colback=white, colframe=aimlsuccess,
  boxrule=0.5pt, left=8pt, right=8pt, top=6pt, bottom=6pt,
  fonttitle=\bfseries\small, coltitle=white, colbacktitle=aimlsuccess,
  attach boxed title to top left={yshift=-2mm, xshift=4mm},
  boxed title style={boxrule=0pt, arc=1pt}, #1}

% ---- small semantics -------------------------------------------------
\newcommand{\scope}[1]{\textcolor{aimlmuted}{\small #1}}
\newcommand{\term}[1]{\textbf{#1}}
\newcommand{\code}[1]{\texttt{\small #1}}
\newcommand{\lecture}[1]{Lecture~#1}

\newtheorem{definition}{Definition}[section]
\newtheorem{proposition}[definition]{Proposition}

% ---- title block -----------------------------------------------------
% \notestitle{19}{Information retrieval: the lexical foundation}{SciFact (BEIR)}
\newcommand{\notestitle}[3]{%
  \thispagestyle{empty}
  \begin{flushleft}
    {\color{aimlmuted}\small\sffamily
     APPLICATIONS OF MACHINE LEARNING \quad$\cdot$\quad
     BSc Mathematics of Artificial Intelligence}\\[2mm]
    {\color{aimlrule}\rule{\textwidth}{0.8pt}}\\[4mm]
    {\color{aimlmuted}\large Lecture #1 \quad$\cdot$\quad Part V \quad$\cdot$\quad
     Extended notes}\\[3mm]
    {\Huge\bfseries\color{aimlprimary}\baselineskip=1.15\baselineskip #2\par}\vspace{4mm}
    {\color{aimlmuted}\large Dataset: #3}\\[3mm]
    {\color{aimlrule}\rule{\textwidth}{0.8pt}}
  \end{flushleft}
  \vspace{2mm}
  \begin{keybox}[title={Why these notes exist}]
    Lectures 19--22 sit outside G\'eron, the course's primary text. For those
    four lectures \emph{these notes are the primary source}, and they are
    examined on exactly the same terms as the textbook parts. Everything in
    the slide deck for this lecture is here, with the arguments written out at
    length rather than compressed onto a slide; the numbers are the ones the
    accompanying notebook prints, and every one of them is a measurement on
    the corpus named above rather than a figure quoted from a paper.
  \end{keybox}
  \vspace{1mm}
  {\small\tableofcontents}
  \vspace{2mm}
  \noindent{\color{aimlrule}\rule{\textwidth}{0.4pt}}
  \vspace{1mm}
}


\begin{document}
\notestitle{19}{Information retrieval: the lexical foundation}{SciFact (BEIR), 5{,}183 abstracts}

\section{The problem}

Eighteen lectures of this course have been about prediction: given an input,
produce a label, a number, or a sequence. Retrieval is not that.

The input is a \term{query} and a \term{corpus}. The output is an
\term{ordering} of the corpus. Nobody wants a yes or a no; they want the
evidence, ranked, so that a person can read the top few and decide.

That change of output changes everything downstream of it. There is no fixed
set of classes, because the corpus changes daily. There is no single right
answer, because several documents may be relevant and their order matters. And
you cannot afford to score the whole corpus with a neural network for every
query, because there are millions of documents and the user is waiting.

\subsection{Why it is not classification}

\begin{center}
\small
\begin{tabular}{lll}
\toprule
 & \textbf{Classification} & \textbf{Retrieval} \\
\midrule
Input             & one document              & a query and a corpus \\
Output            & a label from a fixed set  & an ordering of the corpus \\
Classes           & fixed, known at training  & the corpus, changing daily \\
Right answer      & one                       & several, and their order matters \\
Cost per prediction & one forward pass        & the whole corpus, if you are naive \\
\bottomrule
\end{tabular}
\end{center}

You could phrase retrieval as classification --- one class per document --- and
it would be useless. The classes change every time a document is added, there
are millions of them, and each has a handful of examples.

\subsection{The shape of every search system}

Every production search system has the same three stages.

\begin{enumerate}[leftmargin=*]
  \item \term{Retrieve}: millions $\rightarrow$ hundreds. Must be cheap.
  \item \term{Re-rank}: hundreds $\rightarrow$ tens. Can be expensive.
  \item \term{Present}: tens $\rightarrow$ ten. Diversity, freshness, policy.
\end{enumerate}

This lecture is stage~1, and stage~1 decides the ceiling:

\begin{keybox}[title={The first-stage ceiling}]
A document the first stage never returns \emph{cannot} be recovered by any
later stage, however good that stage is. This is why we will keep quoting
recall@100 for a system whose output is ten results. The first stage is not
judged on what the user sees; it is judged on what it makes possible.
\end{keybox}

Stage~2 is \lecture{20}. Stage~3 is mostly not machine learning and mostly not
in this course.

\subsection{An assumption we import}

Whatever we build, a person reads its output, and that person's attention is
not uniform: the first result is examined far more often than the second,
almost nobody reaches the tenth, and the second page effectively does not
exist.

Nothing in this lecture \emph{measures} that. It is an assumption imported from
outside, and it is flagged as one here because every metric in
Section~\ref{sec:deriv} is an attempt to encode it --- and the differences
between those metrics are differences of opinion about the shape of the decay.

\subsection{The corpus, counted}

\begin{center}
\small
\begin{tabular}{lr}
\toprule
\textbf{Quantity} & \textbf{SciFact} \\
\midrule
Abstracts                                  & 5{,}183 \\
Claims with judgements (test)              & 300 \\
Mean abstract length, in tokens            & 225.2 \\
Median abstract length                     & 213 \\
Vocabulary                                 & 35{,}734 \\
Relevant documents per claim, mean         & 1.13 \\
Claims with exactly one relevant document  & 92.3\% \\
\bottomrule
\end{tabular}
\end{center}

The last row matters more than it looks. With one relevant document per claim
most of the time, recall@10 and the reciprocal rank measure nearly the same
thing. On a corpus where every query had twenty relevant documents, they would
not, and half of Section~\ref{sec:deriv} would read differently.

SciFact is used here for one reason: it is small enough that we can score
\emph{every} document for \emph{every} query, exactly, with no approximate
index in the way. Every number in this lecture and the next is therefore a
property of the method rather than of an index setting nobody inspected.

\subsection{The baseline: return the corpus unranked}

Before any method, ask what \emph{no} method scores. Return every document in
its stored corpus order, the same order for every query:

\begin{center}
\small
\begin{tabular}{lr}
\toprule
\textbf{Metric} & \textbf{Fixed order} \\
\midrule
Recall@10        & 0.0033 \\
Reciprocal rank  & 0.0017 \\
NDCG@10          & 0.0012 \\
\bottomrule
\end{tabular}
\end{center}

Effectively zero --- and that is the point. Unlike accuracy under imbalance in
\lecture{3}, a ranking metric has an honest floor near zero. So ``better than
nothing'' is never an argument in retrieval: the comparison that matters is
always against the \emph{previous method}, and for the next thirty years that
method has been BM25.

\section{From filtering to scoring}

\subsection{First: what is a term?}

Before any weighting there is a decision almost nobody writes up: how the text
becomes a list of things to count. Take a real claim from the corpus,
``0-dimensional biomaterials show inductive properties.'' Lower-case it and
split on non-alphanumeric characters and you get six tokens:

\begin{center}\code{0 \ dimensional \ biomaterials \ show \ inductive \ properties}\end{center}

That is already three decisions, and one of them has done damage:
\code{0-dimensional} has become \code{0} and \code{dimensional}, two terms that
mean nothing apart. Count how many abstracts contain each:

\begin{center}
\small
\begin{tabular}{lr}
\toprule
\textbf{Term} & \textbf{Abstracts containing it} \\
\midrule
\code{0}             & 951 \\
\code{dimensional}   & 72 \\
\code{biomaterials}  & 1 \\
\code{show}          & 1{,}012 \\
\code{inductive}     & 2 \\
\code{properties}    & 211 \\
\bottomrule
\end{tabular}
\end{center}

\code{biomaterials} occurs in one abstract of 5{,}183 and \code{show} in
1{,}012. One of those terms decides the answer and the other decides nothing
--- and no weighting scheme has been chosen yet. The outcome is already largely
determined by the tokeniser.

\subsubsection{What we are deliberately not doing: stemming}

Our tokeniser treats the following as unrelated terms:

\begin{center}
\small
\begin{tabular}{llll}
\toprule
\code{cell} 1{,}837 & \code{cells} 1{,}969 & \code{cellular} 514 & \\
\code{infect} 10 & \code{infects} 7 & \code{infection} 337 & \code{infected} 166 \\
\bottomrule
\end{tabular}
\end{center}

A stemmer would merge each row and pool the evidence: a query for
\code{infect} would then reach the 337 abstracts containing \code{infection}
rather than only 10. It also merges things that should not be merged, and the
net effect on scores is corpus-dependent and usually small. We leave it out
because every unexplained step between the text and the number is somewhere a
result can hide --- and because the mismatch it half-solves has a much better
answer in \lecture{20}.

\subsection{The obvious first idea, and its measurement}

Boolean retrieval returns the documents containing \emph{all} the query terms:
\[ \text{result} \;=\; \bigcap_{t \in q} \text{postings}(t). \]
It is exact, it is fast, and it is how search worked for thirty years. Before
reading on, predict: on a corpus of scientific abstracts, and claims written by
scientists, how many abstracts contain every word of a claim?

\begin{failbox}[title={Boolean retrieval, over all 300 test claims}]
\begin{center}
\small
\begin{tabular}{lr}
Claims matching no document at all & 98.0\% \\
Mean documents matching every term & 0.023 \\
Most any claim matched             & 2 \\
\end{tabular}
\end{center}
98.0\% of claims return an \emph{empty page}. Not a bad ranking --- nothing at
all. And the failure is silent, because an empty result set is
indistinguishable from ``no such document exists''.
\end{failbox}

The conjunction is brittle: one word of the claim absent from an otherwise
perfect abstract, and the abstract is excluded entirely. This is the
observation that reorganises the whole field:

\begin{keybox}[title={Stop filtering; start scoring}]
Every document gets a number. The user gets the largest ones. A missing term
then costs \emph{something} rather than \emph{everything}.
\end{keybox}

\subsection{Why scoring is nevertheless cheap}

Scoring every document sounds worse than filtering, not better. Naively, per
query on this corpus, you would touch 5{,}183 documents and read 1.17~million
tokens --- and this is a tiny corpus. A web index is $10^{10}$ documents and
the user is waiting 200~ms.

But almost every one of those documents shares \emph{no} term with the query
and contributes exactly zero to any sensible score. Reading them is the waste,
and the fix is a data structure rather than a model.

\begin{definition}[Inverted index]
For every term, store the list of documents containing it, together with the
count. One term's entry is called its \term{postings list}.
\end{definition}

\begin{verbatim}
post = defaultdict(list)
for i, doc in enumerate(docs):
    for term, tf in Counter(doc).items():
        post[term].append((i, tf))
\end{verbatim}

Four lines, and it is the entire data structure the field is built on. Scoring
then visits only $\bigcup_{t \in q} \text{postings}(t)$ --- the documents that
can possibly score above zero:

\begin{verbatim}
s = np.zeros(N)
for t in query:
    for i, tf in post.get(t, []):
        s[i] += weight(t, tf, length[i])
\end{verbatim}

\begin{center}
\small
\begin{tabular}{lr}
\toprule
Vocabulary (distinct terms)      & 35{,}734 \\
Postings (term--document pairs)  & 633{,}514 \\
Full matrix, were it dense       & 185 million cells \\
\bottomrule
\end{tabular}
\end{center}

The postings occupy 0.3\% of the dense matrix. Sparsity here is not an
optimisation; it is the reason the field is possible at all. Everything that
remains to be decided is the function \code{weight}.

\section{Term weighting, and BM25}

Start from the obvious idea: a document mentioning \emph{biomaterials} five
times is probably more about biomaterials than one mentioning it once. So score
by how often the query terms occur. Two failures arrive immediately.

\begin{enumerate}[leftmargin=*]
  \item \term{Common words dominate.} Every abstract contains \emph{of} and
        \emph{the}, so a query containing them matches everything equally.
  \item \term{Long documents win.} A 900-word abstract contains more of
        everything than a 100-word one and is not thereby more relevant.
\end{enumerate}

The rest of the method is two corrections, one for each failure, and each has a
parameter --- so each can be measured rather than argued about.

\subsection{Saturation, and the parameter $k_1$}

The first correction: the tenth occurrence of a word should not count ten times
the first.
\[
  \frac{\text{tf}\,(k_1+1)}{\text{tf} + k_1} \;\longrightarrow\; k_1 + 1
  \quad\text{as tf}\to\infty .
\]

\begin{center}
\small
\begin{tabular}{lcccccc}
\toprule
tf     & 1 & 2 & 3 & 5 & 10 & 20 \\
factor & 1.000 & 1.310 & 1.462 & 1.610 & 1.743 & 1.818 \\
\bottomrule
\end{tabular}
\end{center}

From one to two occurrences the factor rises by 0.310; from ten to twenty it
rises by 0.075. The curve is bounded, so no single word can dominate a score by
repetition --- which is, incidentally, the first anti-spam device in search.

$k_1$ sets how fast the saturation bites: small $k_1$ saturates almost
immediately, large $k_1$ approaches raw counting. Measured over 150 claims with
$b$ held at 0.4:

\begin{center}
\small
\begin{tabular}{lcccc}
\toprule
$k_1$    & 0.5 & 0.9 & 1.2 & 2.0 \\
NDCG@10  & 0.6883 & 0.6937 & 0.6944 & 0.6926 \\
\bottomrule
\end{tabular}
\end{center}

The spread across a fourfold change in $k_1$ is 0.0061. That is the honest
finding, and it is not the one you would guess from how much has been written
about tuning this parameter: on abstracts, $k_1$ barely matters. It matters on
long documents, where a term can occur fifty times. \emph{Report the corpus
with the parameter, always.}

\subsection{Inverse document frequency}

The second correction weights a term by how rare it is:
\[
  \mathrm{idf}(t) \;=\; \log\!\left(1 + \frac{N - n_t + 0.5}{n_t + 0.5}\right),
\]
for $N$ documents of which $n_t$ contain $t$.

\begin{center}
\small
\begin{tabular}{lrr}
\toprule
\textbf{Term} & \textbf{Abstracts} & \textbf{idf} \\
\midrule
\code{of}  & 5{,}173 & 0.0020 \\
\code{the} & 5{,}159 & 0.0047 \\
\code{and} & 5{,}158 & 0.0049 \\
\bottomrule
\end{tabular}
\end{center}

\code{of} appears in 5{,}173 of 5{,}183 abstracts and earns a weight of 0.0020.
\emph{No stop-word list was needed.} The weighting derives the stop words from
the corpus it is given.

\begin{keybox}[title={idf replaces the stop-word list}]
Classical systems shipped a hand-written list of words to delete. That list is
language-specific and must be maintained; it is corpus-blind, so in a corpus
about the band \emph{The Who} deleting \emph{the} and \emph{who} destroys the
query; and its decision is binary, with nothing between delete and keep.

A weight of 0.0020 is not deletion, it is near-deletion, and it is computed
\emph{from this corpus}. If a word is common here, it is discounted here, and
nobody maintains anything. This is the same move as learning a representation
instead of engineering features --- arriving thirty years earlier.
\end{keybox}

The particular form above is used rather than the older $\log(N/n_t)$ because
it stays well-behaved for a term appearing in almost every document. Whichever
you use, say which: the two agree closely where $n_t$ is small and diverge
exactly where a stop word lives.

\subsection{Length normalisation, and the parameter $b$}

\[
  \text{tf}' \;=\; \frac{\text{tf}}
  {\text{tf} + k_1\!\left(1 - b + b\,\dfrac{|d|}{\text{avgdl}}\right)}
\]
with $b = 0$ meaning no normalisation, so long documents keep their advantage,
and $b = 1$ meaning length is divided out completely.

\begin{center}
\small
\begin{tabular}{lcccc}
\toprule
$b$      & 0.0 & 0.4 & 0.75 & 1.0 \\
NDCG@10  & 0.6861 & 0.6937 & 0.6917 & 0.7011 \\
\bottomrule
\end{tabular}
\end{center}

The spread is again small, and again that \emph{is} the finding: SciFact
abstracts are of similar length --- 225 tokens on average --- so there is
little here for this correction to correct. On web pages, whose lengths span
three orders of magnitude, there would be.

\subsection{BM25, assembled}

\begin{derivbox}[title={BM25}]
\[
  \text{BM25}(q,d) \;=\; \sum_{t \in q} \mathrm{idf}(t)\,
  \frac{\text{tf}_{t,d}\,(k_1+1)}
       {\text{tf}_{t,d} + k_1\!\left(1-b+b\frac{|d|}{\text{avgdl}}\right)}
\]
\begin{itemize}[leftmargin=*,nosep]
  \item the \term{idf} factor: rare terms count for more;
  \item the \term{saturating tf}: the tenth occurrence adds far less than the
        second, because $k_1$ bounds the growth;
  \item the \term{length term}: $b$ decides how much of the document's length
        is divided out.
\end{itemize}
\end{derivbox}

Every piece is there because a specific failure demanded it. That is why BM25
has survived thirty years of replacement attempts, and it is the baseline every
neural retriever is measured against --- including the one in \lecture{20}.

\subsection{One document, scored term by term}

The claim \emph{``Incidence rates of cervical cancer have decreased.''} against
the 274-token abstract BM25 ranked first for it:

\begin{center}
\small
\begin{tabular}{lrrr}
\toprule
\textbf{Term} & \textbf{tf} & \textbf{idf} & \textbf{contribution} \\
\midrule
\code{cervical}  & 3 & 4.736 & 6.786 \\
\code{incidence} & 5 & 3.238 & 5.145 \\
\code{rates}     & 2 & 2.979 & 3.802 \\
\code{cancer}    & 5 & 1.902 & 3.023 \\
\code{have}      & 3 & 1.008 & 1.445 \\
\code{of}        & 6 & 0.002 & 0.003 \\
\midrule
\textbf{Total}   &   &       & \textbf{20.204} \\
\bottomrule
\end{tabular}
\end{center}

\code{cervical} occurs three times and is worth 6.786. \code{of} occurs six
times --- twice as often --- and is worth 0.003. That ratio is the whole method
working, and this abstract is one of the judged-relevant ones.

\subsection{BM25, measured}

\begin{center}
\small
\begin{tabular}{lrr}
\toprule
\textbf{Metric} & \textbf{Fixed order} & \textbf{BM25} \\
\midrule
Precision@1 & 0.0000 & 0.5433 \\
Recall@10   & 0.0033 & 0.7784 \\
Recall@100  & 0.0117 & 0.8852 \\
MRR         & 0.0017 & 0.6350 \\
NDCG@10     & 0.0012 & 0.6611 \\
\bottomrule
\end{tabular}
\end{center}

The right abstract is in the top ten for 77.8\% of claims and top of the list
for 54.3\%. No training, no labels, no gradient --- a data structure and two
corrections.

Hold on to recall@100: 88.5\%. Almost everything relevant is retrieved
\emph{somewhere} in the first hundred. The gap between @10 and @100 is exactly
what re-ranking exists to close, and it is half of \lecture{20}.

\subsection{Which correction was doing the work?}

An argument is not a measurement. Remove each correction in turn, full test
set:

\begin{center}
\small
\begin{tabular}{lr}
\toprule
\textbf{Scoring function} & \textbf{NDCG@10} \\
\midrule
raw term frequency               & 0.0944 \\
no idf                           & 0.5350 \\
no saturation, no length         & 0.4801 \\
no length normalisation          & 0.6490 \\
BM25                             & 0.6611 \\
\bottomrule
\end{tabular}
\end{center}

\begin{itemize}[leftmargin=*]
  \item \term{idf is most of the method.} Removing it costs 0.1261. Without
        it, common words drown the rare ones that identify the document.
  \item \term{Saturation earns its place.} idf alone on raw counts reaches
        0.4801 against the full method's 0.6611 --- roughly 0.18 of NDCG.
  \item \term{Length normalisation is nearly free here}, worth about 0.012,
        which is consistent with the $b$ sweep and with abstracts being of
        similar length.
\end{itemize}

\begin{keybox}[title={The transferable part}]
Not the ordering --- that is a fact about SciFact. It is that each term of a
formula can be \emph{removed and priced}, and that a component nobody has
priced is a component nobody has justified. The same procedure applies to every
architecture in Part~IV.
\end{keybox}

\begin{trybox}[title={Exercise 19.1}]
Set $b = 0$ in the notebook and re-score; then remove the idf factor entirely
and re-score again. Before running either, write down which of the two you
expect to cost more, and by how much. One of them barely matters on this
corpus, and you should be able to say which and why from the corpus statistics
alone.
\end{trybox}

\section{Derivation: evaluating a ranking}\label{sec:deriv}

\subsection{What we are given}

A derivation needs its inputs named. For one query we have

\begin{itemize}[leftmargin=*,nosep]
  \item a ranked list $d_1, d_2, \ldots, d_n$ --- the system's output, an
        ordering of the corpus;
  \item a set $R$ of relevant documents --- the judgements, made by people, and
        treated as truth for now.
\end{itemize}

Everything that follows is a function of those two objects and nothing else. No
scores, no probabilities, no model --- which is exactly why the same metrics
apply to BM25 here and to a transformer in \lecture{20}.

The one assumption is that $R$ is correct and complete. It is neither, and
Section~\ref{sec:pooling} is about what that costs.

\subsection{A worked example, used throughout}

The claim \emph{``Incidence rates of cervical cancer have decreased.''} has
three relevant abstracts in the corpus. BM25 placed them at ranks 1, 2 and 8:

\begin{center}
\small
\begin{tabular}{cccccccccc}
\toprule
\checkmark & \checkmark & $\cdot$ & $\cdot$ & $\cdot$ & $\cdot$ & $\cdot$ & \checkmark & $\cdot$ & $\cdot$ \\
1 & 2 & 3 & 4 & 5 & 6 & 7 & 8 & 9 & 10 \\
\bottomrule
\end{tabular}
\end{center}

\scope{How this example was chosen: it is the first test claim for which BM25
puts at least three relevant abstracts in the top ten, chosen so that the
arithmetic is visible --- a single hit gives a one-term sum and demonstrates
nothing about a discount. It is therefore better than typical, and the honest
summary of the method remains the mean over all 300 claims.}

\subsection{Precision and recall, at a cutoff}

\[
  \text{P}@k = \frac{|R \cap \text{top-}k|}{k},
  \qquad
  \text{R}@k = \frac{|R \cap \text{top-}k|}{|R|}.
\]

P@$k$ asks: of what I showed you, how much was good. R@$k$ asks: of what was
good, how much did I show you.

On the example with $k=10$, all three relevant abstracts are inside the top
ten, so $\text{R}@10 = 3/3 = 1$: perfect recall. And that is the problem. The
third abstract sits at rank 8, where a user may never look, and recall@10
cannot say so. It counts a \emph{set}; it does not read an \emph{order}.

\subsection{The cutoff is arbitrary}

Every fixed $k$ hides the same thing in a different place. Too small, and a
method that fills ranks 11--20 with the answer scores zero, identically to one
that never finds it. Too large, and rank 1 and rank 100 count the same, so the
metric stops rewarding the ordering it exists to measure.

What we actually want is a number that \emph{decreases smoothly} as a relevant
document moves down, rather than falling off a cliff at one chosen position.
Every metric below is a different answer to that single requirement.

\subsection{Reciprocal rank}

\[ \text{RR} = \frac{1}{\text{rank of the first relevant document}} \]

Our example finds one at rank 1, so $\text{RR} = 1.0000$. Averaged over queries
this is \term{MRR}, and over all 300 claims BM25 scores 0.6350.

\begin{center}
\small
\begin{tabular}{lcccc}
\toprule
First hit at rank & 1 & 2 & 5 & 10 \\
RR & 1.0000 & 0.5000 & 0.2000 & 0.1000 \\
\bottomrule
\end{tabular}
\end{center}

Smooth, bounded, and it moves when the order moves. But it stops reading after
the first hit: our ranking scores 1.0000 whether the other two abstracts sit at
ranks 2 and 8 or nowhere at all. That makes MRR right for ``is there a good
answer near the top?'' and wrong for ``did I find the evidence?''.

\subsection{Average precision}

\[
  \text{AP} = \frac{1}{|R|} \sum_{i \,:\, d_i \in R} \text{P}@i
\]
--- the precision at each rank where a relevant document appears.

\begin{center}
\small
\begin{tabular}{ccc}
\toprule
\textbf{Hit at rank $i$} & \textbf{Relevant so far} & \textbf{P@$i$} \\
\midrule
1 & 1 & 1.0000 \\
2 & 2 & 1.0000 \\
8 & 3 & 0.3750 \\
\bottomrule
\end{tabular}
\end{center}

\[ \text{AP} = \tfrac{1}{3}\left(1.0000 + 1.0000 + 0.3750\right) = 0.7917. \]

Every hit contributes, and one further down contributes less because the
precision above it is worse.

\begin{failbox}[title={Where marks are lost}]
Divide by $|R|$, not by the number of hits actually found. Dividing by the hit
count gives \emph{mean precision}, which is a different and more flattering
quantity, and it is the commonest error in this material.
\end{failbox}

\subsection{Where AP stops, and what replaces it}

AP assumes relevance is \emph{binary}. Often it is not. An abstract that
directly tests a claim and one that mentions it in passing are both
``relevant'', and AP cannot prefer the first. Web-search judgements are
routinely graded: perfect, excellent, good, fair, bad.

So there are now two requirements: a \term{gain} that depends on how relevant a
document is, and a \term{discount} that depends on where it sits. Separating
those two is the whole idea of the last metric.

\subsubsection{Why a logarithm, and not something else}

The discount could be any decreasing function. Three candidates, with what each
implies about a user:

\begin{center}
\small
\begin{tabular}{llll}
\toprule
\textbf{Discount} & \textbf{Rank 1 vs 2} & \textbf{Rank 9 vs 10} & \textbf{Says} \\
\midrule
$1$ (none)         & equal   & equal   & position is irrelevant \\
$1/i$              & halves  & $-10\%$ & position is nearly everything \\
$1/\log_2(i{+}1)$  & $-37\%$ & $-4\%$  & position matters, mildly \\
\bottomrule
\end{tabular}
\end{center}

No derivation picks one; they encode different beliefs. The logarithm is
convention, chosen because it is gentle enough that positions 5 and 10 remain
distinguishable while rank 1 still leads. Convention is not truth: if you
report NDCG you have adopted this curve, and you should be able to say so out
loud.

\subsection{Discounted cumulative gain, and its normalisation}

\[ \text{DCG}@k = \sum_{i=1}^{k} \frac{g_i}{\log_2(i+1)} \]

\begin{center}
\small
\begin{tabular}{crr}
\toprule
\textbf{Hit at rank $i$} & $\log_2(i+1)$ & \textbf{contribution} \\
\midrule
1 & 1.0000 & 1.0000 \\
2 & 1.5850 & 0.6309 \\
8 & 3.1699 & 0.3155 \\
\midrule
\textbf{DCG@10} & & \textbf{1.9464} \\
\bottomrule
\end{tabular}
\end{center}

Is 1.9464 good? The question has no answer, because DCG depends on how many
relevant documents exist: a claim with ten relevant abstracts can reach a far
larger DCG than one with three, however well both are ranked. So divide by the
best this query could possibly have scored --- the ideal ranking, all relevant
documents first.

\begin{center}
\small
\begin{tabular}{cr}
\toprule
\textbf{Ideal rank} & \textbf{contribution} \\
\midrule
1 & 1.0000 \\
2 & 0.6309 \\
3 & 0.5000 \\
\midrule
\textbf{IDCG@10} & \textbf{2.1309} \\
\bottomrule
\end{tabular}
\end{center}

\begin{derivbox}[title={NDCG}]
\[
  \text{NDCG}@k = \frac{\text{DCG}@k}{\text{IDCG}@k}
  = \frac{1.9464}{2.1309} = 0.9134 .
\]
\end{derivbox}

The same ranking, under four metrics:

\begin{center}
\small
\begin{tabular}{lr}
\toprule
\textbf{Metric} & \textbf{Value} \\
\midrule
Recall@10 & 1.0000 \\
RR        & 1.0000 \\
AP        & 0.7917 \\
NDCG@10   & 0.9134 \\
\bottomrule
\end{tabular}
\end{center}

Recall calls this ranking perfect. It is not: one abstract sits at rank 8.
\emph{The metric you report decides what you are allowed to notice.}

\subsection{One relevant document, moved}

The simplest possible ranking: one relevant document, and the only question is
where it sits.

\begin{center}
\small
\begin{tabular}{lcccc}
\toprule
\textbf{Rank} & 1 & 2 & 5 & 10 \\
\midrule
RR        & 1.0000 & 0.5000 & 0.2000 & 0.1000 \\
NDCG@10   & 1.0000 & 0.6309 & 0.3869 & 0.2891 \\
Recall@10 & 1.0000 & 1.0000 & 1.0000 & 1.0000 \\
\bottomrule
\end{tabular}
\end{center}

Recall@10 is constant at 1 across the whole row: it cannot distinguish the best
possible ranking from the worst one that still makes the page. RR falls
fastest; NDCG falls more gently, by construction. And with 92.3\% of our claims
having exactly one relevant abstract, \emph{this row is our corpus, most of the
time}.

\subsection{Which one to report}

\begin{center}
\small
\begin{tabular}{ll}
\toprule
\textbf{Question} & \textbf{Metric} \\
\midrule
Is there a good answer at the top?          & MRR \\
Did I find everything, order aside?         & Recall@$k$ \\
Did I find everything, order counted?       & AP \\
Graded relevance, position counted?         & NDCG@$k$ \\
Will the second stage have anything to work with? & Recall@100 \\
\bottomrule
\end{tabular}
\end{center}

\begin{keybox}[title={The rule}]
Choose the metric from the \emph{question}, before you see any scores.
Choosing it afterwards, from the numbers, is how a paper comes to report
recall@100 for the system that lost on NDCG@10 --- and it is the single
commonest dishonesty in this literature.
\end{keybox}

\subsection{Averaging over queries}

One more choice, usually made silently. Given per-query scores, do you average
the \emph{scores} or pool the \emph{counts}?

\begin{itemize}[leftmargin=*,nosep]
  \item \term{Macro}: mean of the per-query metric. Every query counts equally,
        including one with a single relevant document.
  \item \term{Micro}: pool hits and totals across queries first. Queries with
        many relevant documents dominate.
\end{itemize}

Every number in these notes is macro: the mean over 300 claims of that claim's
own metric. MRR and NDCG are essentially always reported that way; precision
and recall sometimes are not, and the two can differ substantially.

The variance across queries is large and almost never reported. A mean NDCG of
0.6611 contains claims scoring 1 and claims scoring 0.

\subsection{What none of these metrics measure}

\begin{itemize}[leftmargin=*,nosep]
  \item \term{Latency.} A ranking that takes a minute has a real score of zero
        and no metric here notices.
  \item \term{Redundancy.} Ten identical relevant documents score perfectly and
        serve the user once.
  \item \term{The unjudged.} Anything the pool never saw is counted as
        irrelevant --- Section~\ref{sec:pooling}.
  \item \term{Harm.} A confidently ranked wrong answer and a merely absent one
        score identically.
\end{itemize}

A metric is a choice of what to look at, and it is therefore also a choice of
what to ignore.

\begin{trybox}[title={Exercise 19.2}]
Take the relevance pattern $\checkmark\,\cdot\,\checkmark\,\cdot\,\cdot\,\cdot\,\cdot\,\cdot\,\cdot\,\checkmark$
(relevant at ranks 1, 3 and 10, three relevant documents in total) and compute
P@10, R@10, RR, AP and NDCG@10 by hand. Then check each against
\code{rank\_metrics} in the notebook. The written paper has asked for exactly
this in each of the last three years.
\end{trybox}

\section{Where the judgements come from}\label{sec:pooling}

Every metric above takes $R$ as given. Somebody had to decide, for 300 claims
against 5{,}183 abstracts, which pairs are relevant. Judging all of them is
1.55~million decisions for this very small corpus. For a web collection it is
not large --- it is impossible.

\begin{keybox}[title={Pooling}]
Run several systems, take the top $k$ from each, judge the union, and treat
everything unjudged as \emph{not relevant}. That last clause is an assumption,
and it is doing a great deal of work.
\end{keybox}

\begin{failbox}[title={What pooling misses}]
\begin{itemize}[leftmargin=*,nosep]
  \item A document no pooled system ranked highly is unjudged, and therefore
        counted as irrelevant --- whether or not it is.
  \item So a genuinely new method, which finds documents the pool never saw, is
        \emph{penalised for succeeding}.
  \item The pool was built from the systems of its day, so an old collection
        quietly encodes the assumptions of old methods.
\end{itemize}
Reported gains on a mature benchmark are partly gains at \emph{matching the
pool}. It is why a method can win on the benchmark and disappoint in
production, and why a leaderboard difference of a point or two is rarely worth
defending.
\end{failbox}

This is the same shape as the label noise of \lecture{3} and the survivorship
of \lecture{2}: the ground truth is a measurement, made by someone, under
constraints.

\subsection{SciFact specifically}

SciFact's judgements come from the authors of the claims, who cited the
abstracts. They are therefore unusually reliable and unusually \emph{sparse}:
92.3\% of claims have exactly one relevant abstract, mean 1.13. An abstract
that supports a claim but was not cited is unjudged, and scored as a miss.

Which means our 0.6611 is a \emph{lower bound} on what a reader would call the
quality of this ranking, not an estimate of it. Reporting it as ``BM25 achieves
0.6611'' without that sentence is the kind of claim this course exists to stop.

\subsection{Graded relevance changes the answer}

SciFact's judgements are binary. Where they are graded, the gain in DCG is
usually taken as $g_i = 2^{\text{rel}_i} - 1$ rather than $\text{rel}_i$
itself:

\begin{center}
\small
\begin{tabular}{lrr}
\toprule
\textbf{Grade} & \textbf{Linear gain} & $\mathbf{2^{\text{rel}}-1}$ \\
\midrule
0 --- bad       & 0 & 0 \\
1 --- fair      & 1 & 1 \\
2 --- good      & 2 & 3 \\
3 --- excellent & 3 & 7 \\
\bottomrule
\end{tabular}
\end{center}

The exponential says an excellent document at rank 3 is worth more than two
fair ones at ranks 1 and 2 --- a claim about users, not about mathematics, and
one you should be prepared to defend when you make it. Two systems can trade
places when the gain function changes and nothing else does. State which you
used.

\section{The failure BM25 cannot fix}

\begin{failbox}[title={The worst single-evidence claim in the test set}]
\emph{``In mice, P. chabaudi parasites are able to proliferate faster early in
infection when inoculated at lower numbers than when inoculated at high
numbers.''}

BM25 ranks its one relevant abstract \textbf{4{,}999 of 5{,}183}.

\begin{center}
\small
\begin{tabular}{lr}
Query terms in the vocabulary        & 19 \\
Terms the relevant abstract contains & 2 \\
Which terms                          & \code{in}, \code{to} \\
Mean idf of the shared terms         & 0.0378 \\
Mean idf of the missing terms        & 3.2928 \\
\end{tabular}
\end{center}

The abstract shares \code{in} and \code{to} with the claim it is evidence for.
There is nothing for BM25 to match on, and no parameter setting rescues this.
\end{failbox}

That was the worst case. It is not an isolated one. Over every judged
claim--abstract pair in the test set:

\begin{center}
\small
\begin{tabular}{lr}
\toprule
Mean fraction of query terms absent from the relevant abstract & 40.9\% \\
Pairs sharing fewer than half the query terms                  & 28.6\% \\
\bottomrule
\end{tabular}
\end{center}

A relevant document is missing 40.9\% of the query's words on average. People
restate, abbreviate, use the Latin name, or describe the mechanism instead of
naming it --- and a method whose only evidence is shared strings cannot follow
them.

\begin{keybox}[title={The structural limit}]
BM25 matches \emph{strings}. Relevance is about \emph{meaning}. That gap is not
a tuning problem, and it is exactly what an embedding is for --- which is
\lecture{20}, and the reason Part~V comes after Part~IV rather than before it.
\end{keybox}

\subsection{The ceiling, and what is left on the table}

\begin{center}
\small
\begin{tabular}{lrr}
\toprule
\textbf{Metric} & \textbf{BM25} & \textbf{The most it could be} \\
\midrule
Recall@10  & 0.7784 & 1.0000 \\
Recall@100 & 0.8852 & 1.0000 \\
NDCG@10    & 0.6611 & 1.0000 \\
\bottomrule
\end{tabular}
\end{center}

There are two different opportunities here and they need different methods.
Between recall@10 (77.8\%) and recall@100 (88.5\%) sit documents BM25 found but
ranked too low --- a \emph{re-ranking} problem, solvable by an expensive model
over a hundred candidates. The remaining 11.5\% were never retrieved at all; no
re-ranker can reach them, and only a first stage that matches on something
other than words can. Those are the two halves of \lecture{20}.

\section{Reference}

\subsection{Five questions to ask of any retrieval result}

\begin{enumerate}[leftmargin=*]
  \item \term{Which metric, and at what cutoff?} ``Better retrieval'' with no
        metric named is not a claim.
  \item \term{Against which baseline?} If BM25 is absent from the table, ask
        why. It is free, and it is often close.
  \item \term{Whose judgements, produced how?} Pooled from which systems, and
        how long ago?
  \item \term{Same corpus, same queries, same preprocessing?} A different
        tokeniser is a different experiment.
  \item \term{How large is the difference against the spread across queries?}
        A gain of 0.005 on 300 queries is not a finding.
\end{enumerate}

\subsection{Where students lose marks}

\begin{itemize}[leftmargin=*,nosep]
  \item Computing DCG but forgetting to divide by the ideal, and then comparing
        across queries.
  \item Using $\log_2(i)$ rather than $\log_2(i+1)$, which divides by zero at
        rank 1.
  \item Dividing AP by the number of retrieved hits instead of $|R|$ --- that
        is mean precision, not average precision.
  \item Quoting a metric with no cutoff. ``NDCG of 0.66'' is incomplete.
  \item Describing BM25 as ``tf--idf''. It is tf--idf \emph{with the two
        corrections that make it work}.
\end{itemize}

\subsection{Vocabulary}

\begin{center}
\small
\begin{tabular}{ll}
\toprule
Inverted index  & term $\rightarrow$ the documents containing it, with counts \\
Postings list   & one term's entry in that index \\
idf             & a term's weight, falling as more documents contain it \\
BM25            & saturating tf, length-normalised, idf-weighted \\
MRR             & mean of $1/(\text{rank of first hit})$ \\
NDCG@$k$        & discounted gain over the best possible discounted gain \\
qrels           & the relevance judgements themselves \\
Pooling         & judging the union of several systems' top results \\
\bottomrule
\end{tabular}
\end{center}

\subsection{Scope}

\term{Examinable.} The retrieval problem and why it is not classification; the
inverted index and why sparsity makes retrieval possible; idf, saturating tf
and length normalisation, and BM25 assembled from them; Section~\ref{sec:deriv}
in full --- P@$k$, R@$k$, RR, AP, DCG, NDCG and which question each answers;
pooling, and what unjudged documents do to a reported score.

\term{Not examinable --- engineering.} Index compression, skip pointers, query
parsing, and any particular search library.

\term{Beyond the syllabus.} The probabilistic derivation of BM25 from the RSJ
model; query expansion and relevance feedback; learning to rank.

\subsection{Reading}

\begin{itemize}[leftmargin=*,nosep]
  \item These notes are the primary source; the textbook does not cover
        retrieval.
  \item Manning, Raghavan and Sch\"utze, \emph{Introduction to Information
        Retrieval} --- chapters 1, 2 and 6 for the index and term weighting,
        chapter 8 for evaluation. Free online.
  \item Robertson and Zaragoza, \emph{The Probabilistic Relevance Framework:
        BM25 and Beyond} --- where BM25 comes from, if you want the derivation
        we skipped.
  \item The SciFact and BEIR papers, for how these judgements were made.
\end{itemize}

\subsection{Summary}

\begin{keybox}[title={One line to keep}]
A ranking is scored by a question, and the question must be chosen before the
scores are seen.
\end{keybox}

Retrieval returns an ordering of a corpus, not a label, and the metrics,
baselines and methods all follow from that. The inverted index makes it cheap:
postings are 0.3\% of the dense matrix. BM25 is three corrections to counting
words, each answering a stated failure, and it reaches 0.6611 NDCG@10 with no
training at all. A ranking metric answers a question --- recall@10 called our
worked example perfect while an abstract sat at rank 8. And every one of those
numbers rests on judgements someone made on a pool, with everything unjudged
counted as irrelevant.

\end{document}
